Geospatial data underlies some of the most pressing scientific questions and societal issues of our time. Across various scientific domains, this data is growing in volume and complexity, yet supporting its analysis across its various facets remains a significant challenge. This dissertation proposal advances visual analytics as tools for complex geospatial data analysis, addressing three fundamental facets of complexity: (1) multivariate structure, (2) ensemble variability, and (3) multimodality and predictive modeling. Through three major research threads, I aim to develop methods and systems that make multifaceted geospatial data interpretable and actionable.

First, I address the challenge of visualizing multivariate geospatial data expressed across differing spatial discretizations. Despite the ubiquity of such data, effective visualization designs remain underdeveloped. Here, I introduce HexTiles, a domain-agnostic hexagonal-tiling visual encoding that combines semantic icons with hexagonal tessellation to simultaneously encode multiple variables and guide user attention to their interactions. HexTiles explicitly addresses the Modifiable Areal Unit Problem through confidence encodings derived from weighted variance, surfacing within-tile variability that is typically lost in aggregated geospatial visualizations. I apply HexTiles to California water management and Texas presidential election data, demonstrating comparable quantitative performance to existing visual encodings while providing notable improvements in qualitative metrics related to cognitive load and user experience.

Second, I develop methods for analyzing geospatial ensemble data — collections of model runs that are fundamentally similar yet meaningfully distinct. Understanding variability among ensemble members is critical in, for example, climate science in the context of climate change, yet existing analysis techniques leave room for improvement. Here, I introduce ClimateSOM, an interactive visual analytics workflow that combines self-organizing maps and large language models to enable exploration and
interpretation of climate ensemble datasets. I show that the distributional abstraction that the interactive workflow supports can produce insights that may be overlooked with standard data exploration methods. To demonstrate its broader analytical value, in a subsequent and ongoing work, I extend ClimateSOM with statistical hypothesis testing capabilities. I enhance ClimateSOM with statistical rigor, enabling statistically significant detection of distributional shifts in ensemble data. Using this enhanced system, I demonstrate that ClimateSOM can discover and test climatologically novel insights with an investigation of precipitation co-variability patterns within the U.S. West Coast. I show that precipitation relationships are fundamentally structurally reorganizing under climate change, with implications for California water management and climate adaptation. Third, I address the challenge of integrating heterogeneous, multimodal data sources for flood risk prediction. Here, I present DELUGE, a multimodal deep learning framework pluvial flood damage prediction at 1km spatial resolution and daily temporal resolution for the highest-claim areas of Conterminous United States (CONUS), as defined by the National Flood Insurance Program (NFIP) claims.
%
DELUGE explicitly structures prediction around the core disaster-management components of hazard, exposure, vulnerability, and resulting loss.
% DELUGE decomposes flood risk into Hazard, Exposure, and Vulnerability components, integrating geospatial foundation model embeddings to enable lightweight CONUS-scale prediction. 
%
The proposed system outperforms our tabular baselines, and I show that a purely data-driven model can learn pluvial flood damage inducing conditions as a supplement to standard process-based models which often have computational requirements that limit CONUS-scale application. We also introduce a schema to leverage geospatial foundation models with \emph{interpretability-by-design} by routing the embeddings through a physically meaninful parametric modules. This framework establishes the predictive foundation for a visual analytics system for flood risk assessment that I consider future work in this dissertation.

For the remainder of my dissertation, I will continue developing visual analysis methods for complex, multifaceted geospatial data. To consolidate the HexTiles work, I plan a larger and more diverse user study, since the initial evaluation gave strong evidence of reduced cognitive load but, with only seven participants, was underpowered to resolve the quantitative task-performance comparisons. I then pursue three main directions. First, I will build a visual analytics system to accompany DELUGE. Prior work has consistently shown the role that visual analytics can play to bridge the gap between complex predictive models and domain stakeholders, particularly in high-stakes contexts like flood risk assessment. 
% I plan to leverage explainable AI techniques to illuminate
% DELUGE's predictions explicitly along the Hazard, Exposure, and Vulnerability components, fostering interpretability and stakeholder trust. 
Second, I will develop a visual analytics system for understanding how trained models use the embeddings they draw from geospatial foundation models, grounding a model's reliance in human-understandable spatial concepts and surfacing the predictive structure it uses that no known concept explains. Third, I will continue my investigation of climate ensemble data by developing methods to analyze the spatial co-occurrence of extreme climate events which we expect to become more frequent. Standard domain methods, while effective at surfacing point-wise frequency changes, may struggle to capture shifts in how such events co-occur across space. Taken together, this dissertation advances visual analytics as tools for complex geospatial analysis, with particular focus on imminent environmental challenges, equipping analysts with the methods needed to derive interpretable and actionable insights from data.

